\subsubsection{Changing one Hilbert weight}
\label{par:mixed-hilbert}

The two places over the residual prime have different roles in the induction. Changing the Hilbert weight at one while retaining both ordinary lines produces the required local orders.

\begin{proposition}[Simultaneous cuspidal lift at a mixed weight]\label{prop:mixed-hilbert-lift}
For the marked Eisenstein data over the real quadratic field
\(k_0\), the ordinary specialization below lifts the residual
system with one weight changed, retaining the specified unit-root
ordinary lines. Let \(t_w=t\) at the selected place and \(t_w=0\) at the other. The holomorphic weights are \(k+1+2t_w\). The original and lifted eigenvalues obey 
\[
 \ell^{t_w}\lambda^{\prime}_w\equiv\lambda_w\pmod{\mathfrak m}\quad(w\mid\ell),\qquad \lambda^{\prime}_w\equiv\lambda_w\pmod{\mathfrak m}\quad(w\nmid\ell),
\]
 after finite coefficient extension.
\end{proposition}

\begin{proof}
Let \(P\) consist of the place of \(k_0\) above \(\ell\)
at which the Eisenstein unit root is the Hecke value of
\(\eta\). At the other place it is \(1\). Write \(c(w)\)
for the conductor exponent of the Hecke character associated
to \(\eta\) away from \(\ell\), and put \(c(w)=1\) at both
places above \(\ell\). We need a cuspidal system of the
same central character and level with holomorphic weights
\(k+1+2t_w\), where \(t_w=t\) on \(P\) and \(t_w=0\)
at the other embedding. The comparison required of its
eigenvalues is
\[
 \ell^{t_w}\lambda'_w\equiv\lambda_w\pmod{\mathfrak m}
       \quad(w\mid\ell),\qquad
 \lambda'_w\equiv\lambda_w\pmod{\mathfrak m}
       \quad(w\nmid\ell).
\]
Here and below reduction of a normalized eigenvalue uses a
finite extension of \(\Q_\ell\) containing it.

We give a proof using the definite quaternion algebra available
over this real quadratic field. Its two ramified places are
the two real places, so all finite Hecke operators keep their
local double cosets. Global Jacquet--Langlands transfers a
cuspidal representation that is discrete series at both real
places to this algebra. In the opposite direction, an
automorphic representation of the quaternion algebra of
dimension greater than one transfers back with the same finite
components. At a real place the compact type for weight
\(h\ge2\) restricts to \(\operatorname{Sym}^{h-2}\) on the
norm-one group. The central character fixes its extension.
\hfill\mbox{\cite{jacquet-langlands}}

Put \(d_0=k-1\). At the two split \(\ell\)-adic places use
the integral coefficient module
\[
 V_{\mathbf t}=\bigotimes_{w\mid\ell}
       \bigl(\operatorname{Sym}^{d_0+2t_w}(O^2)\otimes\det^{-t_w}\bigr).
\]
For a scalar matrix \(zI\), the two actions on one tensor factor multiply to
\[
 z^{d_0+2t_w}\det(zI)^{-t_w}
 =z^{d_0+2t_w}(z^2)^{-t_w}=z^{d_0}.
\]
Thus the central exponent is independent of the shift at \(w\). Use a sufficiently large finite
coefficient DVR \(O\). We describe the
Hecke normalization. If \(D\) denotes the quaternion algebra,
write \(x\in D^\times(\mathbb A_f)\) and use finite level
\(K_c=\prod_w K_1(w^{c(w)})\) under the finite splittings.
Quaternionic forms are functions invariant on the left by
\(D^\times(k_0)\), of the given central character with its
algebraic normalization, and satisfying
\[
 f(xu)=\rho_{\mathbf t}(u_\ell)^{-1}f(x)
\]
for \(u\in K_c\).
Here \(\rho_{\mathbf t}\) is the displayed coefficient action.
For a double coset with right cosets \(a_iK_c\) the operator
is \(\sum_i\rho_{\mathbf t}(a_{i,\ell})f(xa_i)\).

To see its scalar normalization, start in the unitary
normalization with the compact type at infinity, write
the form as a vector-valued function on the finite ideles,
and transfer the algebraic action on its values to \(\ell\)
using the fixed embeddings. Multiply its value at \(x\) by
\[
 |\operatorname{nrd}x|_f^{-d_0/2}\,
                \rho_{\mathbf t}(x_\ell)^{-1}.
\]
For a rational quaternion the real normalization of its
algebraic coefficient action contains the factor
\(|N_{k_0/\Q}\operatorname{nrd}x|^{-d_0/2}\). The product
formula thus gives the asserted left invariance.

For a finite central idele \(z\), apart from the finite-order
character, the factor is
\[
             |z|_f^{-d_0}\prod_{w\mid\ell}z_w^{-d_0}.
\]
Because \(\ell\) splits in \(k_0\), the valuation of the
central multiplier is
\[
  d_0\sum_{w\mid\ell}v_\ell(z_w)
                       -d_0\sum_{w\mid\ell}v_\ell(z_w)=0 .
\]
Both coefficient modules have this same central factor. The scalar multiplying a
degree-one Hecke double coset relative to its unitary action
is \((Nw)^{d_0/2}\), for both weights. At \(w\mid\ell\) we
further multiply by \(\ell^{t_w}\) to obtain an integral
operator. Indeed its representatives there have the form
\(\left(\begin{smallmatrix}\ell&b\\0&1\end{smallmatrix}\right)\),
\(b\bmod\ell\). For such a representative the normalized matrix on the coefficient module is
\[
 \ell^{t_w}\det(a_i)^{-t_w}\operatorname{Sym}^{d_0+2t_w}(a_i)
 =\operatorname{Sym}^{d_0+2t_w}(a_i) .
\]
It preserves the integral symmetric power because \(a_i\) has integral entries. Away from \(\ell\) the representatives may
be taken to have identity \(\ell\)-components.

The double quotient by the center and \(K_c\) is finite by
finiteness of the ideal classes of a definite quaternion
order and finite index of the level. Evaluation at its
representatives gives a direct sum of coefficient modules
fixed by their projective stabilizers, with the indicated
central factors. The stabilizers are finite: the rational
projective group is discrete adelically, the archimedean
group is compact, and the finite level is compact modulo
the center.

At the ramified place of \(\eta\) of norm two
the positive level forces the local projective stabilizer
into a conjugate of the image of \(K_1(w^{c(w)})\).
Its reduction is contained in the order-two unipotent subgroup
\[
 \left\{\begin{pmatrix}1&a\\0&1\end{pmatrix}:a\in\mathbf F_2\right\},
\]
and its congruence kernel is pro-\(2\). The rational
projective group maps injectively to the local projective
group, since a quaternion which becomes scalar there was
already scalar. For the reduction map on a finite stabilizer \(H_x\), both the kernel and the image have 2-power order. Hence
\[
 |H_x|=|\ker(\mathrm{red}|_{H_x})|\,|\mathrm{red}(H_x)|
\]
is a power of two.
\[
 |H_x|=|\ker(\mathrm{red}|_{H_x})|\,
       |\mathrm{red}(H_x)| .
\]
The kernel is contained in a pro-2 group, and the image lies in the upper unipotent group over \(\mathbf F_2\).

For a projective stabilizer \(H_x\) at a coset \(x\),
the \(O\)-linear projector on its fixed values is
\(e_x=|H_x|^{-1}\sum_{h\in H_x}\rho_x(h)\), including
the fixed central multipliers in the action \(\rho_x\).
The order of \(H_x\) is a power of two and \(\ell\) is odd,
so \(e_x\) is integral. The splitting \(V=e_xV\oplus(1-e_x)V\) reduces to the analogous splitting over the residue field. For an equivariant surjection \(u:V\to W\), if \(w=e_xw=u(v)\), then
\[
 w=e_xu(v)=u(e_x v).
\]
This proves surjectivity on the two fixed submodules. Summing over the finite coset set
proves both assertions for forms.

In each symmetric power let the last-coordinate quotient
be given by the coefficient of the pure power of the
second basis vector. Modulo \(\mathfrak m\) the matrices
of \(K_1(w)\) are upper triangular, so this is a quotient
of their action. For \(t\) a positive multiple of \(\ell-1\)
the determinant factor on the units is trivial after
reduction. The two weights induce the same character on
the last-coordinate quotient, since all diagonal reductions
lie in \(\mathbf F_\ell^\times\). The normalized operators
at \(w\mid\ell\) also act by the same local quotient map:
their representatives have lower right entry one and the determinant factor is
\[
 \ell^{t_w}\det(a_i)^{-t_w}=1.
\]

Tensor the two quotient maps. Write \(Q\) for the product of last-coordinate maps
on the residue modules of forms. In either weight,
\[
                   \ker Q\subseteq\ker(U_{w_1}U_{w_2}) .
\]
This inclusion is at coefficient reduction modulo \(\mathfrak m\)
and uses the two normalized operators. Indeed for a local representative \(a_i\), modulo \(\ell\)
its matrix on the standard space is
\(\begin{pmatrix}0&\bar b\\0&1\end{pmatrix}\).
Write \(D_w=d_0+2t_w\), the exponent in the coefficient
module \(V_{\mathbf t}\). On its symmetric power the
reduced matrix acts by
\[
 e_1^a e_2^{D_w-a}\longmapsto
 \begin{cases}
 0,&a>0,\\
 (\bar b e_1+e_2)^{D_w},&a=0 .
 \end{cases}
\]
The determinant normalization contributes one. Thus after
taking both factors, every tensor monomial in the kernel
of the two last-coordinate quotients is killed. The
level preserves the quotient and its kernel, so the
calculation holds for all double-coset representatives.

Normalize an eigenvector in the starting weight to be
primitive in the integral lattice of forms. Its reduction \(\bar f\) is nonzero and satisfies
\[
 U_{w_1}U_{w_2}\bar f
          =(\bar a_{w_1}\bar a_{w_2})\bar f\ne0
\]
because both eigenvalues are units. The inclusion just proved
therefore gives \(Q(\bar f)\ne0\). The quotient of the mixed-weight lattice surjects onto the
same residual Hecke module. If \(\mathfrak m\) is the
character's maximal ideal, this gives
\[
                    (M_{\mathrm{mixed}})_\mathfrak m\ne0 .
\]
Thus the residual character is in its support.

Let \(\mathbb T\) be the finite torsion-free algebra generated by the commuting operators on the lattice, and choose a minimal prime \(P\) below the residual maximal ideal. A uniformizer \(\pi\) acts injectively on \(\mathbb T\), whereas each element of a minimal prime is a zero divisor. Hence
\[
 \pi\notin P,\qquad (\mathbb T/P)[1/\pi]\ne0 .
\]
In the normalization of \(\mathbb T/P\), choose a prime above
\(\mathfrak m/P\) by lying over. Completing at this prime yields
a characteristic-zero character over a finite coefficient extension
whose reduction is the prescribed residual character. 

This character occurs on the lattice after inverting \(\ell\) and extending
scalars, by faithfulness and simultaneous triangularization.
The degree \(d_0>0\) of the compact type at the unchanged
embedding rules out all one-dimensional quaternionic
representations. Transfer back gives the claimed cuspidal
Hilbert eigensystem. The two normalized eigenvalues above \(\ell\) reduce to the
prescribed nonzero elements \(\bar a_{w_1},\bar a_{w_2}\), so both
ordinary lines survive. The weight can be chosen with
\[
                    t\in(\ell-1)\mathbb Z_{>0},\qquad t\ge2
\]
as large as required for the subsequent local construction.
\label{cor:ordinary-mixed-hilbert-1}

\end{proof}

\subsubsection{Local and global images of the rank-two seed}

An ordinary eigenform with the desired weights supplies a Galois representation. Using that representation after extending the base field requires control of its image on open subgroups.

\begin{proposition}[Absolute irreducibility of the mixed-weight seed on open restrictions]\label{prop:mixed-hilbert-restriction}
The resulting cuspidal rank-two system has the local images used
for the ordinary seed. In the chosen residual situation, with
\(\bar\eta^2\ne1\), let \(V\) be the two-dimensional space of \(r_2\) over \(\overline{\mathbb Q}_\ell\). Absolute irreducibility on open restriction means 
\[
 \forall U\le_{\rm open}G_{k_0},\qquad \{W\subseteq V:\dim W=1,\ r_2(U)W=W\}=\varnothing.
\]

\end{proposition}

\begin{proof}
Apply local--global compatibility for a cuspidal Hilbert
eigenform of discrete-series weights at least two, all of
the same parity and with algebraic central normalization.
The mixed weights just constructed have these properties.
At \(w\nmid\ell\), newform theory and the fixed vector at
\(K_1(w^{c(w)})\) give the upper bound for the local parameter
\(W_w\):
\[
               c(w)=a(\det W_w)\le a(W_w)\le c(w).
\]
The middle inequality follows by comparing invariant
codimensions over each ramification subgroup.

For the finite inertial part of the
Weil--Deligne parameter, the codimension of invariants on
each ramification subgroup is at least that for the
determinant, and monodromy can only increase the
conductor. At positive conductor equality in the inertia
term forces a one-dimensional fixed space and distinct
inertial characters \(1,\eta\). In the two distinct inertia lines the commuting operator \(N\) is diagonal. Consequently
\[
 N=\operatorname{diag}(n_1,n_2),\qquad N^2=0
 \quad\Longrightarrow\quad n_1=n_2=0 .
\] At zero conductor
the parameter is unramified. This gives the full inertia
action away from \(\ell\), in particular at the norm-two
place.
\hfill\mbox{\cite{skinner-hilbert}}\label{cor:ordinary-mixed-hilbert-2}

At the two split places the normalized ordinary eigenvalues
are units. In the normalization before the final cyclotomic
twist the determinant is \(\epsilon_\ell^{-k}\eta\).
The two exponents are
\[
 m_0=t_w,\qquad m_1=-k-t_w,\qquad m_0+m_1=-k.
\]
A nonzero \(U_w\)-eigenvalue
at \(K_1(w)\) level supplies a Borel Jacquet exponent.
With arithmetic Weil degree its normalized Weil character is
\(s\mapsto(\lambda'_w)^{-\deg s}\). Indeed the translation
by the uniformizer on the Jacquet module gives the \(U_w\)
eigenvalue after the Borel modulus factor.

Passing to the
algebraic Weil parameter cancels this factor. The second
Weil character is unramified too since the determinant is
unramified on the Weil inertia group. The Hodge filtration
jumps are \(-m_0,-m_1\). The normalized unit reading and the determinant give the
Newton slopes
\[
        v_\ell(\lambda'_w)=-m_0,\qquad
        k-v_\ell(\lambda'_w)=-m_1 .
\]
Here crystalline Frobenius is inverse to arithmetic Weil
Frobenius in the stated convention. On the
potentially semistable module the smaller-slope line is
stable under monodromy, which lowers slope by one. It is
stable under descent, since each slope occurs once. Its
induced Hodge number is at least the smaller jump. Weak
admissibility gives the reverse bound by the Newton number,
so the line is weakly admissible. The equivalence for
semistable representations and descent give a Galois line.
Its character \(\kappa_0\) has inertia \(\epsilon_\ell^{m_0}\)
by the rank-one comparison.

On an arithmetic Frobenius
\(\sigma\) the same comparison gives
\[
     (\kappa_0\epsilon_\ell^{-m_0})(\sigma)^{-1}
                 = \ell^{m_0}\lambda'_w .
\]
The other character follows from the determinant. The two reductions of
the first character are \(\bar\eta\) on \(P\) and \(1\)
on its complement, by the weight congruence and local
reciprocity. Twisting by \(\epsilon_\ell^k\) leaves the reductions fixed
since \(\ell-1\mid k\). The determinant becomes
\(\epsilon_\ell^k\eta\) and the ordered exponents become
\[
       (m_0+k,m_1+k)=(k+t_w,-t_w).
\]
These are \((k,0)\) and \((k+t,-t)\) at the respective places.
Away from the level and \(\ell\), comparison of the Hecke
polynomials with the Eisenstein system gives residual trace
\(1+\bar\eta\) and determinant \(\bar\eta\). In particular at good Frobenius elements \(g\) the residual polynomial is
\[
 X^2-(1+\bar\eta(g))X+\bar\eta(g)=(X-1)(X-\bar\eta(g)).
\]
Chebotarev and Brauer--Nesbitt identify the residual semisimplification with \(1\oplus\bar\eta\).
\hfill\mbox{\cite{skinner-hilbert,colmez-fontaine}}\label{cor:ordinary-mixed-hilbert-3}

For irreducibility one can also use the unequal gaps at the
two embeddings. If the representation were reducible, the
two constituents would be global \(\ell\)-adic characters,
locally algebraic at both places by the flags and of finite
inertial image away from \(\ell\). In a totally real field
the cyclotomic exponents of such a character are parallel.
Indeed, restrict by reciprocity to totally positive units in
the kernels of all finite inertial factors. A finite-index
subgroup of the totally positive units remains, and the
product of its embedded values to the indicated integer
powers is one.

For the finite-index group of totally positive units,
transfer their algebraic identities to \(\R\). The local
exponents \(a_\sigma\) then satisfy
\[
          \sum_\sigma a_\sigma\log|\sigma(\epsilon)|=0
                        \qquad(\epsilon\text{ in that group}).
\]
Their logarithms span the trace-zero hyperplane by the unit
theorem, so all \(a_\sigma\) are equal. The possible exponents \(k,0\) at one place and
\(k+t,-t\) at the other are disjoint.

Finally let \(H\) be the normal core of any open subgroup.
An inertia element at the unchanged place has infinite
cyclotomic image and a positive power belongs to \(H\).
The two eigenvalues of this power are distinct. If \(H\)
had an invariant line, all such lines would be among
the two eigenlines of that element, and the full group
would permute them. Global irreducibility forces a
transitive action on two lines. The character \(\nu\)
of this permutation is a nontrivial quadratic self-twist:
in the basis of the two lines, conjugation by
\(\operatorname{diag}(1,-1)\) twists the representation by
\(\nu\). Reducing characteristic polynomials gives
\[
 \{1,\bar\eta\}=\{\bar\nu,\bar\nu\bar\eta\}
\]
as
multisets of characters. The nontrivial quadratic \(\nu\) has nontrivial reduction since \(\ell\ne2\). The two entries must therefore be exchanged:
\[
 \bar\nu=\bar\eta,\qquad \bar\nu\bar\eta=1
 \quad\Longrightarrow\quad\bar\eta^2=1,
\]
contrary to the order of the chosen \(\bar\eta\). Thus every open restriction
is absolutely irreducible.
\end{proof}
